% !TeX root = ../main.tex

\chapter{结论}
\section{基础复现保留了什么，尚未复现什么}
本文以申万宏源柏顿·墨基尔成功选股法为对象，保留持续正成长与相对低估值两项筛选主线，并明确初始五年、成熟期历史、全市场排名、正PE和季度持有规则。基础复现B实际覆盖2003年4月至2015年6月147个月，年化38.86\%、累计收益5481.49\%、最大回撤$-$70.11\%。申万原文公布截至2015年7月31日年化36.74\%、累计收益4176.87\%、最大回撤幅度70.19\%。二者在收益数量级与深度回撤上相近，但并非相同日期、相同频率和相同名单的精确重建。

末期原文8只与B的9只仅重合4只，空仓时段也有差别。独立CSMAR诊断能够定位4只未匹配股票的历史同比和报告期缺口，原文发布日与固定季度末的调仓安排亦有明确差异。因此，合理的结论是完成了有明确近似边界的月频规则复现，而不是以年化接近宣布逐日一致。E年化26.17\%包含2015年后的长期经历，不能与原文36.74\%直接相减作为复现误差。

\section{四条路径的收益风险结论}
P1恢复正PE与初始五年基础规则，同时保留成熟期19/20、候选池排名、扩展候选与三维排序。全期年化29.64\%高于E，但最大回撤仍为$-$68.87\%，并未解决极端损失。其改善应归于当前整套规则，不能单独归功于某一个因子或将历史上更高收益等同于更合理的规则。

P2保持其他环节并恢复全截面两榜独立排名，年化27.38\%、最大回撤$-$64.15\%。它相对P1降低了全期年化与最深峰谷损失，却略微提高波动；原复现区间收益排序与全期相反。全市场排名的主要研究价值是恢复相对于基础版可解释的母体，同时揭示排名、补位及最终核心之间的机制，不能笼统称为全面优于P1。

P3、P4分别叠加相同仓位规则，年化降至23.19\%、20.85\%，最大回撤收窄至$-$32.74\%、$-$34.86\%。其风险改善伴随约6.5个百分点的全期年化代价，且零仓位月份均为0。已观察到的作用主要来自波动目标缩仓，不能声称25\%季度止损已经通过真实清仓得到验证；每一阶段的保护程度也不同。

原始时期补测还确认，147、148、151个月三种口径下，P1和P2的收益均高于基础版，P3和P4则低于基础版；151个月中两条选股路径分别提高年化2.33、4.52个百分点。单独考察2008年4月至2015年7月成熟期时，四路径收益均高于B，但受控路径仍低于各自未控底层。这些分组结果要求同时保留收益、风险、起止区间和配对对象，不能笼统声称四条路径在原始时期都提高绝对收益。

\section{持续性、程序证据与适用范围}
初始60个月共同的高复利明显影响全期结果，成熟期及后续阶段需要分开阅读。P1和P2的阶段排序变化，追加21个月中E又领先四路径，表明全期历史最优不是持续优势证明。完整年度计数、滚动五年和回撤恢复记录提供互补描述，均不是未经数据选择的独立样本外检验。

计算方面，B与E在147个共同月逐月一致，四路径初始20个季度256条名单与基础一致，全部正式持仓PE为正。独立基准重建E的560条股票—季度记录及282个月收益，最大月收益误差约$5\times10^{-16}$；R又对17组口径直接复算收益与回撤。上述证据支持既定模型的算术、连接和规则实现一致性，不能证明申万日频等价、原始数据绝对正确或策略具有可交易超额收益。

本研究因此共同保留B、E与P1—P4，而不根据一个终点淘汰其他路径。后续最优先事项是CSMAR历史资料及公告时钟核验、可成交收益和真实成本验证，再以冻结规则评价新增数据。本文的价值在于建立可读、可复核、能解释有利与不利结果的研究链条，并将收益追求、排名定义与风险暴露分别评价。
